# Profil E1 — Manager, 10 ans d'ancienneté

**Fonction :** Manager  
**Ancienneté :** 10 ans  
**Missions :** Banque, assurances, institutions européennes (notamment EIF). Chef de projet sur des livrables techniques et PowerPoints.  
**Rapport à l'IA :** Sceptique — n'utilise pas tant que ça, préfère ses habitudes, se trouve plus rapide que l'IA.  
**Tic de langage :** "À mon niveau"  
**Habitude notable :** Harcèle les juniors sur la relecture et veut savoir exactement ce qui a été généré par l'IA vs écrit à la main.  
**Contexte spécifique :** A vécu une expérience marquante où un livrable junior était truffé d'hallucinations — a dû tout refaire en catastrophe.  
**Outils utilisés :** Copilot de temps en temps, rien d'autre.

---

## T0 — Cadrage (parcours et missions)

> "Manager ici depuis maintenant dix ans. J'ai commencé comme analyste, j'ai gravi les échelons. Je travaille beaucoup avec le secteur bancaire et les assurances, mais aussi avec des institutions européennes — l'EIF notamment, où on produit des livrables techniques assez pointus. Aujourd'hui, je supervise une ou deux missions en parallèle, parfois plus. J'encadre un ou deux juniors en ce moment, mais il y a quelques années, j'en avais plutôt trois ou quatre. La charge a un peu évolué. Mon travail, c'est beaucoup de pilotage — je suis chef de projet sur pas mal de choses. Donc des PowerPoints, des suivis, des revues de livrables. À mon niveau, je passe plus de temps à relire qu'à écrire."

---

## T1 — Usage réel (Appropriation)

### Ouverture — dernière fois que vous avez utilisé un outil génératif

> "La dernière fois... honnêtement, ça doit faire deux ou trois semaines. Et encore, c'était plus par curiosité qu'autre chose. Un junior m'avait parlé d'un prompt qu'il utilisait pour structurer des notes de réunion. Je l'ai testé sur un compte-rendu que j'avais déjà fait à la main. Le résultat était correct, mais il n'apportait rien de plus que ce que j'avais déjà. À mon niveau, je sais ce que je veux, je sais comment le formuler, je n'ai pas besoin d'expliquer pendant deux heures à une machine ce que je veux avant qu'elle comprenne. Je suis plus rapide tout seul."

### Relance — sur quelles tâches au quotidien ?

> "Je ne l'utilise pas vraiment au quotidien. De temps en temps, Copilot, quand je dois faire un premier jet sur un sujet que je maîtrise moins. Mais c'est rare. Je préfère mes méthodes — j'ai des templates, des bases de données d'anciens livrables, une bonne connaissance des attentes des clients. L'IA ne m'apporte pas grand-chose. Par contre, je vois bien que les juniors l'utilisent tout le temps. Eux, ils ont besoin de cette béquille."

### Relance — comment as-tu appris à t'en servir ?

> "J'ai suivi une session de sensibilisation il y a quelques mois, comme tout le monde. Mais franchement, je n'ai pas creusé. Je n'ai pas de bloc-notes avec des prompts, je ne fais pas de veille sur le sujet. Je sais que c'est un outil, je sais à peu près ce qu'il peut faire, mais je ne me suis jamais dit 'il faut que je maîtrise ça'. À mon niveau, ce qui compte, c'est le résultat final, pas le processus pour y arriver."

### Relance — le temps gagné, tu l'aurais mis sur quoi ?

> "Le temps gagné, je ne le vois pas vraiment, parce que je ne l'utilise pas. Mais si je devais imaginer, je dirais que ce gain de temps, il faut le mettre sur la relecture et la vérification. Parce que ce qui sort de l'IA, ça demande du temps à contrôler. Donc au final, est-ce qu'on gagne vraiment du temps ? Je ne suis pas sûr. À mon niveau, je préfère faire une fois, bien, plutôt que de passer du temps à corriger ce qu'une machine a fait approximativement."

### Relance — une situation où ça t'a demandé plus d'effort que de le faire toi-même ?

> "La plus marquante, c'était il y a quelques mois. Un livrable pour un client important, préparé par un junior qui avait utilisé l'IA sur à peu près tout. Je ne m'étais pas plongé dans le projet avant — j'avais une confiance relative dans le junior, et la deadline arrivait. Je reçois le document, je le lis, et là... des incohérences partout. Des chiffres qui ne correspondent pas, des références à des textes qui n'existent pas, des formulations qui semblent justes mais qui trahissent le sens. J'ai dû tout reprendre en catastrophe, avec le junior, pendant une nuit. On a passé plus de temps à corriger ce que l'IA avait fait qu'à produire le livrable nous-mêmes. Depuis, je suis très, très prudent. Et je demande systématiquement : 'qu'est-ce qui est de toi, qu'est-ce qui est de l'IA ?'"

### Relance — vous en parlez entre vous, avec les autres managers ?

> "Oui, pas mal. Entre managers, on partage nos expériences. Il y a une forme de consensus : l'IA est là, il faut composer avec, mais il faut garder un œil très critique. On a tous eu des histoires de livrables qui posaient problème — certains ont même dû être refaits entièrement. À mon niveau, je vois ça comme un changement de métier : je ne vérifie plus seulement le fond et la forme, je vérifie aussi ce que l'IA a apporté. Et ça, c'est une compétence nouvelle que je dois développer."

---

## T2 — Apprentissage du métier (Compétences)

### Ouverture — comment on apprend le métier aujourd'hui ?

> "Quand j'ai commencé, il y a dix ans, on apprenait en faisant, en se trompant, en se faisant reprendre par son manager. Aujourd'hui, je vois des juniors qui arrivent avec des livrables très propres très vite. Mais je me demande parfois si la profondeur est là. L'IA leur donne une première version, ils la retravaillent, mais ils n'ont pas fait le chemin de la réflexion initiale. À mon niveau, je vois des jeunes qui sont très bons sur la forme mais qui ont du mal à justifier le fond. L'apprentissage est différent, et je ne suis pas sûr qu'il soit meilleur."

### Relance — par quelles tâches on apprend réellement le métier ?

> "Les tâches où on est confronté à des choix difficiles. Quand tu dois décider de la recommandation, de l'orientation stratégique, de ce qui est pertinent pour le client. Ça, l'IA ne peut pas le faire. Mais pour les tâches répétitives, la rédaction, la mise en forme, l'IA peut aider — et c'est peut-être une bonne chose. Parce que ça permet aux juniors de passer plus de temps sur les tâches à valeur ajoutée. Dans l'absolu, je ne suis pas contre. Mais il faut que les juniors aient quand même fait les tâches ingrates au moins une ou deux fois à la main, pour comprendre ce qu'il y a derrière. Sans ça, ils ne sauront pas juger ce que l'IA leur sort."

### Relance — comment vous vérifiez ce que produit un junior qui utilise l'IA ?

> "Je ne vérifie plus seulement le contenu final. Je vérifie le processus. Je demande : 'comment es-tu arrivé à cette conclusion ? Qu'est-ce que tu as fait toi-même ? Qu'est-ce que l'IA t'a donné ? Est-ce que tu as vérifié chaque source ?' Je veux savoir ce qui vaut quelque chose ou pas, d'après lui. Parce que ce qui sort de l'IA, ça peut être brillant ou catastrophique, et il faut que le junior sache faire la différence. C'est devenu une partie importante de mon rôle de manager : former les juniors à cette vérification critique."

### Relance — la relecture senior, qu'est-ce qui se transmet à ce moment-là ?

> "À mon niveau, je transmets surtout le jugement. Pourquoi telle formulation est meilleure qu'une autre ? Pourquoi cette recommandation est pertinente et pas celle-là ? Pourquoi ce chiffre est suspect et doit être vérifié ? La relecture, ce n'est pas une correction technique, c'est une transmission de compétences. Avec l'IA, cette transmission devient encore plus importante, parce que les juniors peuvent avoir l'illusion que l'outil fait le travail à leur place. Mon rôle, c'est de leur montrer que non, le jugement reste leur responsabilité."

---

## T3 — Client et valeur (Légitimité de l'expert)

### Ouverture — la confiance du client repose sur quoi ?

> "La confiance du client repose sur notre capacité à être solides. Dans la banque et l'assurance, on ne peut pas se permettre d'être approximatif. Le client doit sentir que ce qu'on lui donne est vérifié, solide, qu'on a pris nos responsabilités. Si l'IA a fait des erreurs et qu'on ne les a pas vues, la confiance est rompue. Et elle ne se reconstruit pas facilement."

### Relance — le sujet de l'IA a-t-il déjà été abordé avec un client ?

> "Oui, plusieurs fois. Certains clients nous demandent si on utilise des outils d'IA. On répond toujours que nous utilisons des technologies d'aide à la décision, mais que l'analyse finale et les recommandations sont entièrement réalisées par nos équipes. C'est la vérité, en fait. Mais on ne dit jamais 'on génère des slides avec Copilot'. Je pense que ce serait mal perçu — le client pourrait se demander ce qu'il paie."

### Relance — diriez-vous à un client qu'une partie a été produite par l'IA ?

> "Non, franchement non. À mon niveau, je pense que ce serait contre-productif. Le client ne paie pas pour le processus, il paie pour le résultat. Et si le résultat est bon, peu importe comment on y est arrivé. Par contre, si le résultat est mauvais, on ne peut pas se cacher derrière l'IA. La responsabilité est entièrement sur nous."

### Relance — si un livrable était perçu comme moins crédible, pourquoi ?

> "À cause d'erreurs factuelles. C'est arrivé — un livrable où un junior avait utilisé l'IA sans vérifier, et il y avait des coquilles grossières. Le client ne nous a pas dit directement, mais on a senti une forme de méfiance après. La prochaine mission, il a demandé un autre consultant. Depuis, je suis intransigeant sur la relecture. À mon niveau, je ne peux pas me permettre de laisser passer une erreur."

### Relance — dans dix ans, sur quoi un client acceptera-t-il de payer un cabinet ?

> "Il paiera pour le jugement et la responsabilité. L'IA pourra faire de plus en plus de choses, mais elle ne pourra pas prendre de risque à la place du consultant. Dans la banque, dans l'assurance, la prise de risque est centrale. Le client paie quelqu'un qui dit 'je vous recommande ça, et j'assume'. L'IA ne peut pas faire ça. Donc la valeur du cabinet, ce sera le jugement — pas la production."

---

## T4 — Gouvernance et risques

### Ouverture — il existe des règles sur l'usage des outils ici ?

> "Oui, il y a une politique. Copilot est l'outil officiel, on a des règles sur la confidentialité — ne pas mettre de données sensibles, etc. Mais à mon niveau, je trouve que les règles sont un peu légères. Il y a un flou sur les outils non-officiels, et surtout, il n'y a pas de processus de vérification spécifique. Ça repose sur la responsabilité individuelle. Ce qui est bien, mais pas suffisant quand on voit les risques qu'on a déjà rencontrés."

### Relance — avant qu'un livrable ne parte chez le client, il se passe quoi ?

> "Il y a une relecture, c'est moi ou un autre manager. Mais jusqu'à récemment, on ne vérifiait pas spécifiquement ce qui venait de l'IA. Depuis l'incident dont j'ai parlé, je vérifie. Je demande aux juniors de me montrer leurs sources, de me justifier leurs choix. C'est devenu un réflexe de sécurité. À mon niveau, je pense que les managers doivent développer cette compétence de vérification spécifique — parce que l'IA change la nature des erreurs."

### Relance — déjà rencontré un contenu produit par l'IA qui s'est révélé faux ?

> "Oui, comme je l'ai dit — un livrable entier. Des chiffres faux, des références inventées, une structure qui ne tenait pas la route. Heureusement, on a pu le refaire avant que le client le voie. Mais on a passé une nuit à le reprendre, et on a livré avec quelques heures de retard. Le junior a eu une leçon, et moi aussi. Depuis, j'ai changé ma manière de superviser. À mon niveau, je ne peux pas avoir confiance aveugle dans ce que font les juniors avec l'IA — je dois tout contrôler."

### Relance — si vous écriviez la règle vous-même, vous diriez quoi ?

> "Je dirais : 'L'IA est un outil d'assistance, pas de production. Chaque livrable doit pouvoir être justifié par le consultant sans recours à l'outil.' Et j'ajouterais une obligation de traçabilité — un système qui permet de savoir ce qui a été produit par l'IA et ce qui ne l'a pas été. Pas pour contrôler les gens, mais pour pouvoir vérifier en cas de problème. À mon niveau, c'est le minimum pour éviter les catastrophes."

---

## T5 — Clôture

### Relance — à la place de la direction, quelle serait votre première décision ?

> "Je mettrais en place un système de validation renforcé pour les livrables qui utilisent l'IA. Pas une interdiction, mais un marquage — un peu comme on fait pour les documents confidentiels. Une mention qui indique 'ce livrable a été préparé avec l'assistance d'outils d'IA' et qui déclenche une relecture renforcée. Et je formerais les managers à cette relecture spécifique. Parce que les erreurs de l'IA sont différentes des erreurs humaines — elles sont plus subtiles, plus difficiles à détecter."

### Relance — un aspect important qu'on n'a pas abordé ?

> "La pression commerciale. Les clients demandent de plus en plus d'IA — ils veulent des consultants qui maîtrisent ces outils. Le cabinet doit s'adapter, et les managers comme moi doivent montrer l'exemple. Mais je trouve que c'est un peu contradictoire : on nous dit de maîtriser l'IA, mais on ne nous donne pas de formation approfondie, pas de temps pour apprendre. À mon niveau, je dois à la fois garder mes compétences techniques et apprendre ces nouveaux outils. C'est une pression supplémentaire."

### Relance — quelqu'un dont le point de vue serait particulièrement utile ?

> "Un directeur ou un associé. Pour comprendre comment ils voient l'évolution du métier et comment ils intègrent l'IA dans la stratégie du cabinet. Moi, je vois les risques au niveau des livrables. Eux, ils voient les risques au niveau de la marque et de la réputation. Et ils doivent aussi gérer la pression des clients qui demandent des solutions IA. Leur regard serait très différent du mien."